\chapter*{Abstract}
\addcontentsline{toc}{chapter}{Abstract}

\begin{spacing}{1.4}

Accurate fetal biometry depends on the localisation of anatomical
landmarks on standard ultrasound planes, from which measurements of
fetal growth are derived. Manual landmark placement is subject to
inter-observer variability, motivating automated landmark localisation.
Existing approaches are largely based on task-specific convolutional
architectures, whereas recent self-supervised vision foundation models
offer general-purpose visual representations learned through
large-scale pretraining. Yet adapting such representations to precise
anatomical landmark predictions without sacrificing the spatial
accuracy required for fetal biometry remains underexplored.

This thesis introduces \textbf{QCLand} (Query-Conditioned Landmark
Localisation), a framework for adapting pretrained vision foundation
models to fetal landmark localisation. For each measurement-specific
model, QCLand incorporates two learned landmark query tokens into a
pretrained vision foundation model, allowing the queries to interact
with image tokens and form landmark-specific representations through
the transformer's self-attention. A new spatial decoding head,
DeconvHeadV2, then converts these query-conditioned representations
into spatial landmark heatmaps for precise localisation. The framework
is instantiated with two self-supervised backbones, DINOv2 and DINOv3,
yielding QCLand-DINOv2 and QCLand-DINOv3.

QCLand is evaluated on five fetal biometric measurements across the
UCL and Multicentre datasets using permutation-invariant normalised
mean error (PI-NME), and is compared with reproduced HRNet-W18 and
RTMPose-s baselines. QCLand-DINOv3 achieves the lowest five-seed mean
PI-NME on all five UCL measurements. Across the two datasets and five
measurements, a QCLand variant records the lowest mean error in eight
of the ten comparisons. On Multicentre, a QCLand variant records the
lowest mean PI-NME on both abdominal measurements and femur length,
while HRNet-W18 records the lowest mean on the two head measurements.
The BPD development experiments further show that introducing
DeconvHeadV2 produces the largest observed reduction in localisation
error among the evaluated architectural modifications.

These results provide evidence that vision foundation models can be
effectively adapted for precise fetal landmark localisation when
combined with landmark-query conditioning and an appropriate spatial
decoder. QCLand emerges as a competitive alternative to purpose-built
convolutional landmark-localisation architectures, supporting
query-conditioned foundation-model adaptation as a viable route to
precise anatomical landmark localisation.

\end{spacing}
